\chapter{CertOS Runtime and Certificate Lifecycle}
\label{ch:certos-runtime}

This chapter promotes CertOS from roadmap language into a bounded runtime
surface backed by code and locked artifacts.  The thesis claim remains
intentionally narrow.  CertOS is \emph{not} presented here as a universal
runtime proof or a deployment announcement.  It is presented as a battery
reference runtime that enforces certificate-aware execution, maintains an
append-only audit chain, and falls back safely when certificate validity is
exhausted.

\section{Why the Runtime Layer Exists}

The DC3S battery shield prevents unsafe actions at the optimisation
boundary.  A deployment-facing system still needs a second question
answered: \emph{what happens between a safe action and the runtime that is
allowed to execute it?}  CertOS answers that question with one simple
discipline:

\begin{quote}
\textbf{No action without a valid certificate or an explicit fallback
action.}
\end{quote}

This runtime discipline matters for three reasons.  First, it makes
certificate validity a live operational object rather than a static chapter
concept.  Second, it turns every intervention into an auditable lifecycle
event rather than an implicit side effect.  Third, it gives the battery
thesis a controlled path from batch evidence toward controller-in-the-loop
execution without pretending that the full deployment problem has already
been solved.

\section{CertOS Code Surface}

The implementation lives under \path{src/orius/certos/}.  The main runtime
file is \path{runtime.py}; the surrounding modules handle lifecycle
issuance, audit logging, recovery, reachability, reliability, and safe
action filtering.  Table~\ref{tab:certos-module-map} maps the main runtime
surface to the concrete code paths.

\begin{table}[ht]
\centering
\caption{CertOS runtime module map and bounded role in the battery thesis.}
\label{tab:certos-module-map}
\begin{tabular}{p{3.2cm}p{4.8cm}p{4.6cm}}
\toprule
\textbf{Module} & \textbf{Primary role} & \textbf{Code path} \\
\midrule
Runtime orchestrator & Per-step lifecycle dispatch and invariant checks & \path{src/orius/certos/runtime.py} \\
Certificate engine & ISSUE, VALIDATE, EXPIRE, RENEW, REVOKE, FALLBACK & \path{src/orius/certos/certificate_engine.py} \\
Audit ledger & Append-only runtime event log & \path{src/orius/certos/audit_ledger.py} \\
Recovery manager & Recovery callback after expiry/fallback & \path{src/orius/certos/recovery_manager.py} \\
Belief engine & Combined state and uncertainty object & \path{src/orius/certos/belief_engine.py} \\
Reliability engine & Runtime quality scoring for degraded observation & \path{src/orius/certos/reliability_engine.py} \\
Shift engine & Uncertainty-set construction & \path{src/orius/certos/shift_engine.py} \\
Reachability engine & Validity-horizon computation & \path{src/orius/certos/reachability_engine.py} \\
Safe-action filter & Final safe action enforcement & \path{src/orius/certos/safe_action_filter.py} \\
Graceful planner & Structured fallback planning & \path{src/orius/certos/graceful_planner.py} \\
\bottomrule
\end{tabular}
\end{table}

\section{Lifecycle Operations}

The certificate engine implements six lifecycle operations.  These are not
just conceptual labels; each one corresponds to a callable path in
\path{certificate_engine.py} and is exercised through
\path{scripts/run_certos_lifecycle.py} and \path{tests/test_certos.py}.

\begin{table}[ht]
\centering
\caption{CertOS lifecycle operations and their runtime meaning.}
\label{tab:certos-lifecycle-ops}
\begin{tabular}{p{1.7cm}p{4.2cm}p{5.0cm}}
\toprule
\textbf{Op} & \textbf{When it occurs} & \textbf{Battery-runtime meaning} \\
\midrule
ISSUE & New safe action and positive validity horizon & Normal certified dispatch step \\
VALIDATE & Check current certificate before action release & Guard against stale or exhausted authorisation \\
EXPIRE & Validity horizon reaches zero & Runtime may no longer trust the previously issued action \\
RENEW & Horizon is still positive but below degraded threshold & Continue operating with explicit degraded semantics \\
REVOKE & Runtime or operator invalidates current certificate & Force transition away from normal execution \\
FALLBACK & No valid certificate remains & Replace actuation with a known-safe fallback action \\
\bottomrule
\end{tabular}
\end{table}

\begin{figure}[ht]
\centering
\begin{tikzpicture}[
  node distance=1.9cm and 2.1cm,
  every node/.style={draw, rounded corners, align=center, minimum width=2.8cm, minimum height=0.9cm},
  >=Latex
]
\node (issue) {ISSUE\\valid};
\node[right=of issue] (renew) {RENEW\\degraded};
\node[below=of renew] (fallback) {FALLBACK};
\node[left=of fallback] (expire) {EXPIRE};
\node[below=of issue] (revoke) {REVOKE};

\draw[->] (issue) -- node[above]{low horizon} (renew);
\draw[->] (renew) -- node[right]{horizon $\le 0$} (fallback);
\draw[->] (issue) -- node[left]{manual invalidate} (revoke);
\draw[->] (revoke) -- (fallback);
\draw[->] (issue) -- node[left]{horizon $\le 0$} (expire);
\draw[->] (expire) -- (fallback);
\draw[->, bend left=18] (fallback) to node[below]{recovery + new cert} (issue);
\draw[->, bend left=18] (renew) to node[below]{recovered quality} (issue);
\end{tikzpicture}
\caption{CertOS lifecycle state machine as implemented in the battery runtime.
The defended contract is simple: remain in valid or degraded certified
execution when a horizon exists; otherwise fall back explicitly.}
\label{fig:certos-lifecycle}
\end{figure}

\section{Runtime Invariants}

The runtime is small enough that its main safety contract can be stated as
three explicit invariants.  These are codified in
\path{CertOSRuntime.check_invariants()} and exercised in
\path{tests/test_certos.py}.

\begin{table}[ht]
\centering
\caption{CertOS runtime invariants.}
\label{tab:certos-invariants}
\begin{tabular}{p{1.6cm}p{4.0cm}p{5.4cm}}
\toprule
\textbf{Invariant} & \textbf{Statement} & \textbf{Operational interpretation} \\
\midrule
INV-1 & No dispatch without a valid or fallback certificate & The runtime must never emit a free-running actuation path \\
INV-2 & Certificate hash chain is unbroken & Every lifecycle event remains auditable and ordered \\
INV-3 & Fallback triggers iff validity horizon is exhausted & The fallback path is explicit rather than discretionary \\
\bottomrule
\end{tabular}
\end{table}

The value of these invariants in the battery thesis is not theoretical
novelty.  Their value is that they turn runtime behavior into something
that can be traced, tested, and audited after the shield has already
produced a safe action.

\section{Battery Runtime Evidence}

The publication-facing runtime artifacts live under \path{reports/certos/}.
They include a 96-step lifecycle log, a JSON summary, a latency report, an
audit-completeness report, and a short recovery note.  Table
\ref{tab:certos-runtime-summary} summarises the main locked battery runtime
surface.

\begin{table}[ht]
\centering
\caption{Locked CertOS runtime summary from the battery lifecycle demo.}
\label{tab:certos-runtime-summary}
\begin{tabular}{lr}
\toprule
\textbf{Metric} & \textbf{Value} \\
\midrule
Total steps & 96 \\
Valid steps & 47 \\
Degraded steps & 16 \\
Fallback steps & 33 \\
Runtime interventions & 33 \\
Audit entries & 96 \\
Audit completeness & 100.0\% \\
Step latency & 0.59\,ms \\
Throughput & 1691.70 steps/s \\
\bottomrule
\end{tabular}
\end{table}

Two facts are especially important.  First, the audit-completeness report
records 96 audit entries for 96 runtime steps, so the lifecycle trace is
complete for the governed demo.  Second, the measured 0.59\,ms step latency
shows that the certificate-aware runtime layer is cheap relative to hourly
battery dispatch timescales.  This does \emph{not} prove suitability for
every real-time control loop; it does show that the current battery runtime
prototype is not bottlenecked by certificate bookkeeping.

\section{Representative Runtime Transitions}

Table~\ref{tab:certos-representative-rows} shows representative lifecycle
states from the battery trace.  The goal is not to report every row in the
main chapter; the full log is preserved in Appendix~\ref{app:certos-lifecycle-log}.
The main chapter only needs enough rows to make the runtime logic legible.

\begin{table}[ht]
\centering
\caption{Representative CertOS lifecycle rows from the governed battery
runtime log.}
\label{tab:certos-representative-rows}
\begin{tabular}{rrrrrl}
\toprule
\textbf{Step} & \textbf{Horizon} & \textbf{Status} & \textbf{Discharge} &
\textbf{SOC} & \textbf{Fallback} \\
\midrule
0  & 18 & valid    & 45.0 & 88.16 & no \\
7  & 20 & valid    & 50.0 & 0.00  & no \\
40 & 0  & fallback & 0.0  & 0.00  & yes \\
56 & 6  & valid    & 15.0 & 0.00  & no \\
57 & 4  & degraded & 10.0 & 0.00  & no \\
73 & 0  & fallback & 0.0  & 0.00  & yes \\
83 & 2  & degraded & 5.0  & 0.00  & no \\
95 & 0  & fallback & 0.0  & 0.00  & yes \\
\bottomrule
\end{tabular}
\end{table}

The pattern is operationally coherent.  Healthy steps stay in a certified
\texttt{valid} regime.  As the horizon compresses, the runtime can remain
active but only in an explicitly marked \texttt{degraded} regime.  When the
horizon reaches zero, the runtime does not improvise; it falls back to a
zero-action battery-safe command.  That is the systems contribution of the
chapter.

\begin{figure}[ht]
\centering
\begin{tikzpicture}[
  node distance=1.8cm and 1.8cm,
  every node/.style={draw, rounded corners, align=center, minimum width=2.4cm, minimum height=0.85cm},
  >=Latex
]
\node (obs) {Observed state\\and safe action};
\node[right=of obs] (validate) {Validate\\certificate};
\node[right=of validate] (act) {Release\\action};
\node[below=of validate] (fb) {Fallback\\action};
\node[right=of act] (audit) {Audit\\ledger};

\draw[->] (obs) -- (validate);
\draw[->] (validate) -- node[above]{valid / degraded} (act);
\draw[->] (validate) -- node[right]{expired / revoked} (fb);
\draw[->] (act) -- (audit);
\draw[->] (fb) -- (audit);
\end{tikzpicture}
\caption{No-action-without-certificate runtime path.  A release path exists
only for certified actions; otherwise the runtime records a fallback event
and emits the fallback action instead.}
\label{fig:certos-runtime-path}
\end{figure}

\section{Recovery, Audit, and Bounded Portability}

The recovery note in \path{reports/certos/recovery_report.md} records that
the lifecycle contains transitions from fallback back to valid execution.
The recovery manager in \path{recovery_manager.py} is deliberately simple:
it is a callback-based regain path, not a proof of optimal recovery.  That
level of modesty is appropriate to the current evidence lock.

The audit layer is similarly bounded.  The ledger stores lifecycle
operation, step, certificate hash, and optional metadata.  The current
battery artifact proves completeness for the demo and preserves an
append-only runtime trail.  It does not yet prove cryptographic immutability
or fleet-scale storage guarantees.

CertOS also includes a small second-domain toy artifact in
\path{reports/certos/certos_vehicle_toy.json}.  That artifact records a
12-step vehicle-shaped run with four fallback steps, twelve audit entries,
and four interventions.  The role of the toy is narrow: it shows that the
certificate lifecycle is not hard-coded to battery field names.  It does
\emph{not} elevate the thesis into a second fully validated proof domain.

\begin{table}[ht]
\centering
\caption{Bounded non-battery portability evidence for the CertOS runtime.}
\label{tab:certos-vehicle-toy}
\begin{tabular}{lr}
\toprule
\textbf{Vehicle toy metric} & \textbf{Value} \\
\midrule
Domain & vehicle \\
Steps & 12 \\
Fallback steps & 4 \\
Audit entries & 12 \\
Interventions & 4 \\
\bottomrule
\end{tabular}
\end{table}

\section{Bounded Claim}

The runtime chapter therefore supports five bounded statements:

\begin{enumerate}
  \item the repo contains a concrete certificate-aware runtime rather than
    only a conceptual deployment story,
  \item the six lifecycle operations are implemented and test-covered,
  \item the battery lifecycle demo records complete audit coverage and
    explicit fallback behavior,
  \item the prototype runtime is computationally lightweight at the measured
    operating scale, and
  \item the lifecycle contract ports to a toy second domain without changing
    the battery thesis into a universal runtime proof.
\end{enumerate}

\section{Summary}

CertOS closes the gap between a safe battery action and a runtime that is
allowed to execute it.  The current chapter does not claim deployment
completion.  It claims something narrower and stronger: the thesis now has a
bounded certificate lifecycle chapter with code anchors, runtime invariants,
measured battery evidence, complete audit coverage, explicit fallback
semantics, and a small portability demonstration that remains outside the
main proof boundary.
